\section{Bilan et livrables}

\subsection{Livrables attendus}

\begin{longtable}{|p{5cm}|p{4cm}|p{4cm}|}
\hline
\textbf{Livrable} & \textbf{Statut} & \textbf{Emplacement / remarque} \\
\hline
Cahier des charges & Réalisé & Présent dans ce dossier LaTeX. \\
\hline
Analyse fonctionnelle & Réalisé & Présente dans ce dossier LaTeX. \\
\hline
Outils de management & Réalisé & WBS, planning, Kanban, risques. \\
\hline
Documentation technique & Réalisé & Architecture, modules, matériel, tests. \\
\hline
Poster de présentation & Hors document & Réalisé séparément. \\
\hline
Archives code & Réalisé & Dossier \texttt{src/systus}. \\
\hline
Archives CAO & Réalisé & \texttt{boitier-systus.stl}, \texttt{support-systus.stl}, \texttt{ecran-systus.stl}. \\
\hline
Schémas électroniques & Partiellement documenté & Schéma fourni et décrit ; image/source à ajouter au dépôt. \\
\hline
Prototype fonctionnel & À valider en soutenance & Démonstration le 11 juin 2026. \\
\hline
\end{longtable}

\subsection{Fichiers PDF à rendre}

\begin{longtable}{|p{5cm}|p{5cm}|p{3cm}|}
\hline
\textbf{Fichier attendu} & \textbf{Contenu} & \textbf{Statut} \\
\hline
\texttt{SYSTUS\_CDC.PDF} & Cahier des charges et analyse fonctionnelle. & Couvert par ce dossier \\
\hline
\texttt{SYSTUS\_DocTechnique.PDF} & Documentation technique, tests, architecture et archives. & Couvert par ce dossier \\
\hline
\texttt{SYSTUS\_Poster.PDF} & Poster de présentation. & Réalisé séparément \\
\hline
\end{longtable}

\subsection{État du prototype au 10 juin 2026}

Au moment du rendu documentaire, le dépôt contient une base logicielle cohérente
pour le prototype :

\begin{itemize}
    \item boucle principale avec modes \texttt{setup} et \texttt{running} ;
    \item capture audio par \texttt{arecord} ;
    \item client de reconnaissance AudD ;
    \item machine à états pour la détection ;
    \item affichage e-paper des états et résultats ;
    \item interface Flask de configuration Wi-Fi ;
    \item gestion de boutons physiques via GPIO ;
    \item fichiers STL du boîtier, du support et de la pièce écran archivés.
\end{itemize}

Les éléments à compléter concernent principalement la preuve matérielle finale :
photos, image ou source du schéma électrique et résultats mesurés pendant les
tests.

\subsection{Conclusion}

S.Y.S.T.U.S. valide l'idée d'un objet musical autonome capable de transformer
une interaction physique simple en reconnaissance musicale et affichage e-paper.
Le projet respecte l'esprit du prototypage : partir d'un besoin, cadrer les
fonctions, construire rapidement une chaîne testable, puis identifier les
limites et les pistes d'amélioration.

La suite logique du projet consiste à sécuriser la configuration, finaliser le
montage matériel, enrichir le boîtier et améliorer la robustesse de la
démonstration.
